% Transcription — fonds Alexandre Grothendieck, shelfmark 150, batch 3 (pages 41-60).
% Established from the facsimile archives/batches/150/batch-03.pdf (a copy of
% 150.pdf fetched through the project's relay,
% https://grothendieck-relay.onrender.com/source/150.pdf; PDF page k =
% archivists' page 40+k, no cover sheet), per the skill transcribe-grothendieck.
% The folder is 98 pages in five batches (1-20, 21-40, 41-60, 61-80, 81-98);
% this is the third, done in parallel with the others and aligned afterwards
% on batches 1 and 2 (transcripts/150/batch-01.fr.tex, batch-02.fr.tex).
%
% Pass: Opus 5.5 (claude-opus-5-5), 2026-09-26 — first pass, unchecked against the pages by a human.
%
% Pages skipped: none. Every page carries his mathematics.
%
% Pages summarised: none. The cancelled lower half of page 57 (four long
% diagonals, a first attempt at (37)-(38)) is transcribed as \struck{}, and
% the struck head of page 46 likewise.
%
% His pagination and numbered runs. The batch continues the run headed
% « Espaces stratifiés (étude par voisinages coniques) » (batch 1, sheets
% 1)-5); batch 2, sheets 6-16). Sheets numbered top left: 17 (p. 42),
% 18 (44), 19 (46), 20 (48), 21 (50), 22 (52), 23 (54), 24 (56), 25 (58,
% in braces), 26 (60); the unnumbered pages 41, 43, 45, 47, 49, 51, 53, 55,
% 57, 59 are their continuations (probably versos). Page 41 continues
% sheet 16 (p. 40), whose last sentence « La donnée des Sigma_ijk ->
% Sigma_ik revient : » runs on here; its margin refers to « p. 15' », his
% sheet 15 (p. 38). Paragraphs: « 10. » (p. 42, circled (10) in the margin,
% no title) and « 11. Digression ... » (p. 47, underlined title, mostly
% illegible). Numbered formulas (24)-(39), continuing batch 2's (23): (24)
% (read also as 29), (25), (25*), (26), (27), (28), (29), (30) (only as a
% reference), (31)-(35), (36) twice (p. 54 and p. 57), (37)-(38) struck on
% p. 57 and restated on p. 58, (39) twice (p. 58 and p. 60) — kept as
% written, with \note{}s. Conditions: circled (a)-(f) on pp. 43-45
% (a « cylindrique », b, c, d transversality, e « équisingularité »,
% f « non-singularité »), again circled (a)-(e) on pp. 47-50 for the
% ordered topological set Sigma, b_1)-b_3) on p. 51, and (a)-(c) once more
% in the restart of p. 55. No date on these pages.
%
% What the batch is about. Pages 41-46: the « présentation standard » of an
% equisingular stratification of type I (begun on pp. 38-40) recast as an
% M-category over I_M; the reduction of the cartesian condition to
% Sigma_ijk = Sigma_ik ∩ Sigma_jk, i.e. (Sigma_*)_{<=z} -> I_{<=k} injective
% with induced order; the simultaneous mapping cylinder C_{Sigma,d,d'} of a
% chain of initial segments; paragraph 10 restates everything for
% M = (Top): Sigma_* with a closed order graph Sigma_**, an increasing map
% Sigma_* -> I, conditions (a) the target maps Sigma_ij -> Sigma_j are closed
% embeddings, (b) (28), (c) disjointness when i, j have no lower bound,
% (d) transversality of the Sigma_ik in Sigma_k, (e) Sigma_ij -> Sigma_i
% fibrations, (f) the Sigma_i manifolds with boundary; the boundary of the
% fibres of Sigma_d' -> Sigma_d stratified by Drap(I_{i_p<.<i_p+1} ⊔ ...);
% smoothness iff each term is a successor of the previous. Pages 47-60:
% paragraph 11, a « digression » without the index set I: a topological
% ordered set Sigma = Sigma(0), Sigma(1) (graph of the order), with the
% semi-simplicial space Sigma*(n) of strict flags; conditions (a) proper
% locally bordant target map, (b) cartesian squares (31)-(32), transversality
% (33)-(34), (c) equisingularity, (d) non-singularity, (e) mixed; the
% reworking b_1)-b_3) in terms of germs B_alpha of bordant subspaces;
% the functor Gamma : d -> Sigma*(d ⊔ {omega}) and the description as an
% « ordered local cylindrical stratification »; Gamma'(d) ≃ S_d × Gamma(d)
% (35); a relative total order Gamma_1 with transitivity (36); then a fresh
% start (p. 55, « Je m'aperçois finalement que j'ai pas mal bafouillé »):
% the order recovered from the source map s = p_1 : Sigma*(1) -> Sigma, the
% map p_2 : Sigma*(2) -> Sigma*(1), the cartesian square (38), the relation
% x ≺ y (39), its transitivity and irreflexivity, and a « Scholie »
% summarising the data (1°-3°), breaking off at 3° a).
%
% Notation, following batches 1-2: Sigma_*, Sigma_**, Sigma_d for d a flag,
% Drap(I), I_M, \underline b, \underline s (and \underline b', \underline s')
% as he underlines them; M as \mathcal M. From page 47 he writes the spaces
% of strict flags with a small raised star, set Sigma^{*}(n) (said once in
% a \note); \mathfrak S_d for the symmetric group in (35); « ½ simpliciales »
% keeps his ½ for « semi- ». His underlinings in running text are \emph{};
% circled numbers and letters are (n) with a \note. Whether « coniques » on
% p. 41 carries his circumflex could not be decided; it is left plain.
%
% Hardest pages: 41, 44-45, 48 and 50 (long sideways notes up the left
% margin, read only in fragments), 46 (struck and looped passage), 57
% (cancelled half). The formulas and numbered conditions carry; much of the
% sideways marginalia does not and is left \ill{}.

\documentclass[11pt,a4paper]{article}
% Le préambule commun à toutes les transcriptions du fonds Grothendieck.
%
% Il tient en une idée : **l'incertitude se compose, elle ne se résout pas.**
% Une transcription qui lisse un mot douteux a détruit la seule chose pour
% laquelle on la faisait — un lecteur doit pouvoir savoir, à chaque endroit,
% ce qui était sur le feuillet et ce qui a été ajouté. Les six macros qui
% suivent sont donc l'appareil critique, et non de la décoration.
%
% Les mêmes commandes sont reconnues par `scripts/render.mjs`, qui produit la
% vue de lecture du site. Une macro ajoutée ici doit l'être aussi là-bas,
% faute de quoi le rendu échouera bruyamment — ce qui est voulu : un rendu
% silencieusement faux est pire qu'un rendu absent.

\ProvidesPackage{grothendieck}[2026/08/08 Transcriptions du fonds Grothendieck]

\usepackage{amsmath,amssymb,amsthm}
\usepackage{mathrsfs}
\usepackage{stmaryrd}
% Les diagrammes commutatifs sont partout dans ce fonds : c'est la langue
% même de la géométrie algébrique de Grothendieck, et une transcription qui
% les rendrait en prose ne transcrirait rien.
\usepackage{tikz-cd}
% Français par défaut — c'est la langue des feuillets — mais l'anglais est
% chargé aussi : la traduction et le résumé sont des documents anglais, et la
% typographie française (espaces avant les deux-points) y serait une faute.
% Ces documents basculent par \selectlanguage{english} en tête de corps.
\usepackage[english,french]{babel}
\usepackage{fontspec}
\usepackage{geometry}
\geometry{a4paper,margin=25mm,marginparwidth=28mm}
\usepackage{marginnote}
\usepackage{xcolor}
% ulem, not soul. `soul`'s \st and \ul are built on a character-by-character
% scan that XeLaTeX's font machinery breaks: the document fails with an
% undefined control sequence rather than degrading. ulem does the same two jobs
% and is Unicode-engine safe. `normalem` stops it hijacking \emph, which the
% transcriptions use for Grothendieck's own emphasis.
\usepackage[normalem]{ulem}
\usepackage{hyperref}
\hypersetup{colorlinks=true,linkcolor=black,urlcolor=black,pdfborder={0 0 0}}

% --- Métadonnées du lot -----------------------------------------------------
% Reprises telles quelles par le script de rendu, qui les affiche en tête de
% la vue de lecture. Elles fixent ce que le fichier prétend couvrir : sans
% elles, une transcription flotte sans point d'attache dans un fonds de seize
% mille feuillets.

\newcommand{\folder}[1]{\gdef\grothFolder{#1}}
\newcommand{\batch}[1]{\gdef\grothBatch{#1}}
\newcommand{\leaves}[2]{\gdef\grothFirst{#1}\gdef\grothLast{#2}}
\newcommand{\foldertitle}[1]{\gdef\grothFoldertitle{#1}}
\gdef\grothFolder{?}\gdef\grothBatch{?}\gdef\grothFirst{?}\gdef\grothLast{?}\gdef\grothFoldertitle{}

% --- Le résumé --------------------------------------------------------------
%
% Il ouvre la lecture modernisée, sous le titre « Résumé », et il est composé
% en petit. Ce n'est pas un ornement : c'est le seul passage du document qui
% ne soit pas des mathématiques, et un lecteur doit pouvoir le reconnaître
% avant de l'avoir lu — puis le sauter s'il connaît déjà le sujet, ou s'y
% arrêter s'il ne le connaît pas. Le filet qui le suit marque la reprise du
% corps.

\newenvironment{resume}
  {\small\setlength{\parskip}{0.45em}}
  {\par\medskip\hrule\medskip}

% --- Mots-clés ---------------------------------------------------------------
%
% En anglais, en clôture du résumé de la lecture modernisée : le vocabulaire
% moderne sous lequel chercher ce que les feuillets construisent. Le manifeste
% du site (`npm run manifest`) les extrait pour étiqueter la cote entière —
% cette ligne est la source unique des tags, il n'y a pas de fichier à côté.
\newcommand{\keywords}[1]{\par\smallskip\noindent
  {\footnotesize\textsc{Keywords} --- \textit{#1}}\par}

% --- L'appareil critique ----------------------------------------------------

% Le numéro de feuillet, en marge. C'est l'ancre qui permet de régler un
% désaccord sur une formule en regardant *un* feuillet plutôt qu'une chemise
% de six cents. Le site s'en sert aussi pour tourner les pages du fac-similé
% quand on fait défiler la transcription.
\newcommand{\leaf}[1]{%
  \marginnote{\footnotesize\color{gray!70}\textsf{#1}}%
  \label{leaf:#1}\ignorespaces}

% Illisible : on ne devine pas. Un mot inventé qui se lit comme les autres est
% le pire résultat possible.
\newcommand{\ill}{\textcolor{red!60!black}{[\,\ldots\,]}}

% Lecture douteuse : le mot est donné, et signalé comme incertain.
\newcommand{\uncertain}[1]{\uline{#1}}

% Ajout de l'éditeur — une lettre manquante, un mot sous-entendu.
\newcommand{\add}[1]{\textcolor{blue!50!black}{[#1]}}

% Biffé par Grothendieck lui-même. Ce qu'il a rayé fait partie du document :
% une rature dit souvent où la pensée a changé de direction.
\newcommand{\struck}[1]{\sout{#1}}

% Note du transcripteur, sur l'état matériel du feuillet.
\newcommand{\note}[1]{\par\smallskip{\small\color{gray!60!black}\textit{[#1]}}\par\smallskip}

% Note marginale de Grothendieck, distincte du corps du texte.
\newcommand{\marginal}[1]{\par\smallskip\noindent\hspace{1em}%
  {\small\color{gray!60!black}#1}\par\smallskip}

% --- Titre ------------------------------------------------------------------

\AtBeginDocument{%
  \begin{center}
    {\large\bfseries Fonds Alexandre Grothendieck --- cote n\textsuperscript{o}~\grothFolder}\\[2pt]
    {\itshape\grothFoldertitle}\\[4pt]
    {\small feuillets \grothFirst--\grothLast\ (lot \grothBatch)}\\[2pt]
    {\footnotesize\color{gray!60!black}Transcription établie d'après le fac-similé numérisé
     par l'Université de Montpellier.\\
     Les lectures incertaines sont soulignées, les illisibles notées [\,\ldots\,],
     les ajouts entre crochets.}
  \end{center}
  \vspace{1em}}

\endinput


\folder{150}
\batch{3}
\pages{41}{60}
\dating{1981-1982}
\watermark{Édition de démonstration}
\foldertitle{Espaces stratifiés et voisinages côniques (ou : déploiement des espaces stratifiés) : notes manuscrites (s.d.)}

\begin{document}

\page{41}
\note{Sans numéro de sa main : suite du feuillet 16 (page 40), dont la dernière phrase, « La donnée des $\Sigma_{ijk}\to\Sigma_{ik}$ revient : », se poursuit ici. Le « (3) » renvoie à la condition (3) de la page 39. Écriture rapide ; la prose de liaison se lit mal, les formules passent.}

une \emph{application} \emph{composition}, et l'axiome : formules pour avoir une extension \add{des $\Sigma_i\Sigma_j$} \struck{\ill{}} en des $\Sigma_{i_0\cdots i_n}$, et que cette composition (qui est compatible avec les unités, \uncertain{manifestement}) \add{\uncertain{par} \ill{}} \uncertain{construction}) soit une structure de \emph{$\mathcal{M}$-catégorie}* \struck{\ill{}}. Donc la donnée de $\Sigma$ \add{\uncertain{satisfaisant} (\uncertain{voir} p.\ 15')} (3) équivaut : celle d'une \add{\uncertain{Catégorie dans}} $\mathcal{M}$, au-dessus de la catégorie $I_{\mathcal{M}}$ :
\[
\Sigma_{\mathcal{M}}\longrightarrow I_{\mathcal{M}} .
\]

\marginal{* \struck{en} \uncertain{en fait}, c'est \uncertain{une catégorie} \emph{\uncertain{ordonnée}} de $\mathcal{M}$, et $\Sigma_{i_0\cdots i_n}$ s'identifie \uncertain{à l'ensemble} des $x_{i_0}\le x_{i_1}\le\cdots\le x_{i_n}$ de $\Sigma_{*}$, au-dessus des drapeaux $i_0,\ldots,i_n$ de $I$. Donc pour vérifier \uncertain{(3)} \ill{}
\[
\Sigma_{i_0\cdots i_n}=\Sigma_{i_0i_1}\cap\cdots\cap\Sigma_{i_0i_n}\ \ill{}
\]
il suffit de le vérifier pour $n=2$, i.e.
\[
\Sigma_{ijk}=\Sigma_{ik}\cap\Sigma_{jk},
\]
ce qui signifie que \ill{} \ill{} \ill{} $\forall z\in\Sigma_k$, $\Sigma_{*}\to I$ induit
\[
\Sigma_{*}(z)\hookrightarrow I(k),\qquad \Sigma_{*}(z)=\lbrace x\in\Sigma_{*}\mid x\le z\rbrace,\quad I(k)=\lbrace i\in I\mid i\le k\rbrace,
\]
ce qui est \uncertain{tel} que l'on \ill{} \ill{} $\Sigma_{*}(z)$ \ill{} \emph{\uncertain{induit}} \uncertain{par celui de} $I(k)$. \struck{\ill{}}}\note{Longue note en travers dans la marge gauche, montant en diagonale, appelée par l'astérisque après « $\mathcal{M}$-catégorie » ; lecture d'ensemble incertaine, les formules sont sûres. Le renvoi « p.\ 15' », ajouté en travers devant « (3) », vise son feuillet 15 (page 38), où sont posées les conditions (1), (2), (3).}

À l'aide de ces données, je \uncertain{vais} construire un espace stratifié de type $I$, \add{\uncertain{(je} n'utilise en fait \add{\uncertain{les conditions}} \struck{la \ill{} des} \struck{hypothèses} (1)} \struck{de (3)}, la condition (2) n'intervenant que sous la forme plus faible que
\[
\Sigma_{i_0\cdots i_n}\longrightarrow\Sigma_{i_0\cdots\hat\imath_p\cdots i_n}\qquad(p\neq n)
\]
est un plongement. \struck{\ill{}} (\uncertain{ceci en vue des} constructions \add{\uncertain{relatives}} \uncertain{de} structures coniques \uncertain{d'espaces stratifiés}\ill{})\note{la marque au-dessus de l'indice $i_p$ est lue comme un chapeau (indice omis).}

Je considère pour ceci les couples $d,d'$ ($d\subset d'$) de deux drapeaux \struck{\ill{}} tels que $d$ \struck{\ill{}} soit cofinal dans $d'$. (Pour $d'$ \add{$=\lbrace i_0,\ldots,i_n\rbrace$}, \struck{de longueur $n$} il y a donc $2^n$ choix possibles pour $d$).

\page{42}
\note{En haut à gauche, « 17 » de sa main : numéro de feuillet de sa suite.}

On considère pour cette donnée équivalente : celle d'une suite d'inclusions de drapeaux
\[
d_1\subset d_2\subset\cdots\subset d_p\subset d_{p+1}=d'\qquad(p=\mathrm{long.}\,d')
\]
dont chacun est un segment initial strict du suivant. \marginal{\uncertain{$d=$} \ill{} $\lbrace i_0,\ldots\rbrace$ \ill{} $d_1,\ldots,d_p$}Ceci donne donc \struck{\ill{}} une suite de projections
\[
\Sigma_{d'}=\Sigma_{d_{p+1}}\longrightarrow\Sigma_{d_p}\longrightarrow\cdots\longrightarrow\Sigma_{d_2}\longrightarrow\Sigma_{d_1}
\]
dont on considère le mapping-cylindre simultané que je désigne par $C_\Sigma(d_1,d_2,\ldots,d_{p+1})$ \add{($d_{p+1}=d'$)}, ou par
\[
C_{\Sigma,d,d'} .
\]
On veut définir les \add{\uncertain{morphismes de}} \uncertain{transitions} naturels entre \uncertain{les} $C_{\Sigma,d,d'}$.\marginal{\uncertain{On} \ill{} \uncertain{construction} \ill{} \ill{} \uncertain{finalement} \ldots}\note{note en travers dans la marge gauche, en partie soulignée ; lecture très incertaine.}

\emph{10.} Reprenons la description \uncertain{\ill{}} (dans le cas $\mathcal{M}=(\mathrm{Top})$, pour fixer les idées -- c'était le \uncertain{contexte adopté au début de ces notes}).

(10)\note{numéro cerclé, dans la marge.} \struck{\ill{}} \uncertain{Alors} on a un espace topologique
\[
(24)\qquad \Sigma_{*},
\]
muni d'une relation d'ordre $\Sigma_{**}\subset\Sigma_{*}\times\Sigma_{*}$ \uncertain{ayant} un graphe fermé, et tel que \add{$\mathrm{pr}_1=\mathrm{but}$} $\Sigma_{**}\to\Sigma_{*}$ soit un \struck{\uncertain{plongement}} \add{\uncertain{\ill{}}} -- plus\note{le numéro de la formule, encadré, se lit aussi bien 29 que 24 ; la suite (25), (26), (27) de la page 43 fait lire 24.}

\page{43}
\note{Suite de la page 42, sans numéro de sa main.}

précisément, on a de plus une application croissante et continue
\[
(25)\qquad \Sigma_{*}\longrightarrow I
\]
\note{dans la marge, devant le numéro, « (2°) ».}
\[
\Big(\text{d'où}\quad \Sigma_{*}=\coprod_{i\in I}\Sigma_i,\qquad \Sigma_{**}=\coprod_{i\le j}\Sigma_{ij}\Big)
\]
\note{sous la seconde égalité, une ligne biffée : \struck{$\Sigma_{ii}=\Sigma_i$}.}

\marginal{condition (a)}\note{« a » cerclé ; le passage qui suit est marqué d'un trait vertical.} et on \struck{\ill{}} suppose que $\forall i,j$ \add{($i\le j$)}, l'application but
\[
\Sigma_{ij}\longrightarrow\Sigma_j\qquad\text{« cylindrique »}
\]
est un \struck{\uncertain{immersion}} plongement fermé (-- ce qui implique (par \uncertain{injectivité}) que (25*) a la propriété que si $z\in\Sigma_{*}$ est au-dessus de $j$ et si $i\le j$, alors il existe \emph{au plus} un $y\in\Sigma_{*}$ au-dessus de $i$, avec $y\le z$.\marginal{$\Sigma_{*}^{\le z}\to I^{\le j}$ \uncertain{injectif}}

Cela implique aussi :
\[
(26)\qquad \Sigma_{ii}=\Sigma_i .
\]

\marginal{\uncertain{condition} (b)}\note{le passage est marqué d'un trait vertical.} Car \struck{\ill{} de plus, que pour $z\in$} \struck{(27) $i<j<k$,} $\Sigma(n)$ = ens.\ des drapeaux stricts de longueur $n$ de $\Sigma$ \add{\struck{$\Sigma_i$}}

et
\[
(27)\qquad
\begin{array}{ccc}
\Sigma(n) & \longrightarrow & \mathrm{Drap}_n(I)\quad(\text{drapeaux stricts})\\
\Vert & & \\
\displaystyle\coprod_{\substack{(i_0<\cdots<i_n)\ \text{drapeau strict}\\ \text{de } I \text{ de longueur } n}}\Sigma_{i_0\cdots i_n} & &
\end{array}
\qquad\text{OK}
\]
\marginal{\uncertain{revenir} \ill{} \uncertain{à la} condition (\ill{})}\note{devant « $\mathrm{Drap}_n(I)$ », un $\Sigma$ biffé ; « OK » est écrit à gauche de la réunion disjointe.}

\page{44}
\note{En haut à gauche, « 18 » de sa main.}

\marginal{condition (b)}\note{« b » cerclé.} On suppose, pour $i<j<k$,
\[
(28)\qquad \Sigma_{ijk}=\Sigma_{ik}\cap\Sigma_{jk}\qquad(\text{a priori } \Sigma_{ijk}\subset\Sigma_{ik}\cap\Sigma_{jk})
\]
ce qui signifie aussi que $\forall z\in\Sigma_{*}$ au-dessus de $k\in I$ (i.e.\ $z\in\Sigma_k$), l'\uncertain{hom.\ déf.}
\[
\Sigma_{*}\longrightarrow I\quad(\text{\uncertain{qui} \add{\uncertain{est} \ill{}} induit une application } \Sigma_{*\le z}\hookrightarrow I_{\le k})\ \text{induit}
\]
\[
(29)\qquad (\Sigma_{*})_{\le z}\hookrightarrow I_{\le k}
\]
l'ordre de $(\Sigma_{*})_{\le z}$ est \emph{induit} par celui de $I_{\le k}$.

\marginal{\uncertain{Considérons} \uncertain{que} si $d,d'$ sont des drapeaux \uncertain{tels} que $d''=d\cup d'$ soit un drapeau et que $d,d'$ \ill{} \ill{}, i.e.\ $d,d'\subset$ \ill{}, alors $\Sigma_{d''}=\Sigma_d\cap\Sigma_{d'}$ dans $\Sigma_i$ \ldots\ \emph{injection}}\note{note en travers dans la marge gauche, à hauteur de (28)--(29) ; lecture incertaine.}

On suppose de plus que si $i,j\le k$ \add{dans $I$} n'ont pas de minorant, alors
\[
\Sigma_{ik}\cap\Sigma_{jk}=\emptyset
\]
\marginal{condition (c)}\note{« c » cerclé ; à côté, « \uncertain{Gros} ».} (si $i,j$ \add{\uncertain{ont}} \uncertain{un minimum} $l$, alors
\[
\Sigma_{ik}\cap\Sigma_{jk}\subset\ \ill{}\ )
\]
\marginal{Implique que si $d,d'$ \ill{} deux drapeaux \ill{} \ill{} et si \ill{} \ill{} \ill{}, alors $\Sigma_d\cap\Sigma_{d'}=\emptyset$}\note{note en travers dans la marge gauche ; lecture incertaine.}

\marginal{(b) et (c) impliquent que, pour une famille $(d_\alpha)$ de drapeaux ayant \uncertain{même} \uncertain{point final} $k$, \ill{} les $\Sigma_{d_\alpha}\subset\Sigma_k$, $\bigcap\Sigma_{d_\alpha}$ est \emph{vide} si $\bigcup d_\alpha$ \uncertain{n'est pas un drapeau}, et $=\Sigma_{\bigcup d_\alpha}$ sinon}\note{note en travers dans la marge droite, « b » cerclé ; lecture incertaine.}

On suppose enfin que le système des $\Sigma_{ik}$ (pour $k$ fixé, \add{$i\in I_{<k}$}), qui sont des sous-espaces \struck{bordants} bordants de $\Sigma_k$, sont \emph{transversaux}, \add{et (28),} ce qui signifie (\uncertain{grâce à} : (30)) que si \uncertain{$d$ est} un drapeau $d=(i_0,\ldots,i_n,k)$, alors les $\Sigma_{i_0,k},\Sigma_{i_1,k},\ldots,\Sigma_{i_n,k}$ dans $\Sigma_k$ sont des sous-espaces bordants qui se coupent \emph{transversalement} le long de $\Sigma_d=\Sigma_{i_0\cdots i_n,k}$.

\marginal{condition (d)}\note{« d » cerclé ; en travers dans la marge, une note : « implique que \ill{} les $\Sigma_d$ \ill{} \ill{} \ill{} $d'$ \ill{} \ill{} \ill{} $\Sigma_{d'}\subset\Sigma_d$ \uncertain{bordant} dans $\Sigma_d$ », lecture très incertaine.}

\page{45}
\note{Suite de la page 44, sans numéro de sa main.}

Enfin, dans le cas \add{\uncertain{d'un espace}} \add{\uncertain{stratifié}} équisingulier, on suppose que les
\[
\Sigma_{ij}\longrightarrow\Sigma_i
\]
\add{$\Sigma_{d'}\to\Sigma_d$ \ill{} \ill{}} \note{ajout interlinéaire en partie barré.}sont des \emph{fibrations} (i.e.\ $\Sigma_{**}\xrightarrow{\ s\ }\Sigma_{*}$ est une fibration), et dans le cas des strates \struck{\ill{}} non singulières, on suppose que les $\Sigma_i$ sont des variétés à bord (i.e.\ $\Sigma_{*}$ variété à bord) -- ce qui implique que les $\Sigma_{ji}$ \struck{($j<i$)} \uncertain{toujours distinctes}) sont des sous-variétés à bord du \struck{$\Sigma_i$} bord $\dot\Sigma_i$ de $\Sigma_i$, et les $\Sigma_{ji}$ ($j<i$) \add{\uncertain{le long de} \ill{} \uncertain{codim.}\ 1} \uncertain{forment} une stratification cylindrique des variétés à bord -- \uncertain{plus généralement} les $\Sigma_d=\Sigma_{i_1\cdots i_n}$ sont alors des variétés à bord et les $\Sigma_{d'}$ ($d'\supsetneq d$, $d$ cofinal dans $d'$)* \struck{\ill{}} forment une stratification cylindrique du bord $\dot\Sigma_d$ de $\Sigma_d$.

\marginal{condition (e) « équisingularité ». Cela implique que \ill{} $\Sigma_d\to\Sigma_{\ill{}}$ \ill{} \ill{} fibrations, si $d\subset d'$ \ill{}}\marginal{condition (f) « non-singularité »}\note{notes en travers dans la marge gauche, en face des deux hypothèses.}

\marginal{* et \uncertain{l'ensemble} des $d'$ s'identifie \uncertain{à} \ill{} \ill{} $\widehat{\mathrm{Drap}}(I_{<i_1})\times\widehat{\mathrm{Drap}}(I_{i_1<\cdot<i_2})\times\cdots\times\widehat{\mathrm{Drap}}(\ill{})$, où $\widehat{\mathrm{Drap}}$ désigne les \ill{} \ill{} \ill{} \ill{} (\uncertain{y compris} \ill{} \ill{}) \ldots\ il s'identifie \uncertain{à} $\mathrm{Drap}(I_{<i_1}\amalg I_{i_1<\cdot<i_2}\amalg\cdots\amalg I_{i_{n-1}<\cdot<i_n})$}\note{note en travers dans la marge gauche, reliée par une flèche à l'astérisque du texte ; lecture incertaine hors des formules.}

Dans le cas équisingulier-non singulier, on \uncertain{suppose} \add{\uncertain{de plus}} (e) et (f), que les fibres des \add{\uncertain{fibrations}}
\[
\Sigma_{ij}\longrightarrow\Sigma_i
\]
sont des variétés à bord -- le \add{\uncertain{fibré sur} $\Sigma_i$} \struck{\ill{}} des bords des fibres est \add{\uncertain{réunion}} \struck{\ill{}} stratifié par les $\Sigma_d$, où $d$ \struck{\ill{}} parcourt les drapeaux qui contiennent \emph{strictement} $\lbrace i,j\rbrace$ et ont $i$ comme premier \struck{élément ; \ill{} produit $\mathrm{Drap}(I_{i<\cdot<j})\times\mathrm{Drap}(I_{>j})$}\note{sous les deux facteurs biffés, les légendes « drapeaux \uncertain{cot.\ stricts} » et « drapeaux \ill{} », également biffées ; la phrase continue page 46.},

\page{46}
\note{En haut à gauche, « 19 » de sa main. Les six premières lignes sont biffées de trois longs traits obliques ; elles continuent le passage biffé de la fin de la page 45.}

\struck{\emph{privé} de $(\emptyset,\emptyset)$ (qui correspond au drapeau $\lbrace i,j\rbrace$ lui-même -- si on ne l'excluait pas, on aurait une stratification de $\Sigma_{(i,j)}$ lui-même, non de son bord. Pour que $\Sigma_{ij}\to\Sigma_i$ soit lisse, il f.\ et s.\ donc que $j$ soit un successeur de $i$, et que $j$ \uncertain{soit maximal dans} $I$}\note{dans la marge, en face de « lisse », « (\uncertain{hyp.\ C.\ I}) » souligné, également sous les traits.}

élément (pour être fibré sur $\Sigma_i$) et $j$ comme dernier (pour être contenus dans $\Sigma_j$) d'où $\simeq\mathrm{Drap}(I_{i<\cdot<j})$. Il est vide ssi $j$ est un successeur de $i$, qui est donc la condition pour que $\Sigma_{ij}\to\Sigma_i$ soit lisse.

Plus généralement, si $d\subsetneq d'$ est un segment initial, la fibration $\Sigma_{d'}\to\Sigma_d$ est à fibres des variétés à bord, et le fibré des bords est stratifié cylindriquement par les $\Sigma_{d''}$, où $d''\supsetneq d'$ (pour que $\Sigma_{d''}\to\Sigma_{d'}$) est tel que $d'$ soit cofinal dans $d''$ (pour que $\Sigma_{d''}\hookrightarrow\Sigma_{d'}$) et que $d$ soit segment initial de $d''$ (pour que $\Sigma_{d''}\to\Sigma_d$ soit fibrant) i.e.\ \struck{si $d$ est le} \add{si $d'=(i_1,\ldots,i_n)$,} $d=(i_1,\ldots,i_p)$, \struck{donc $d''$ est de la forme} d'où $d'=d\amalg d'_0$ \struck{plus grand élément \ill{} de $d$ \ill{} $d'_0=(i_{p+1},i_{p+2},\ldots,i_n)$, alors les $d''$ correspondant aux $d'$ ($d_\alpha$, $\alpha\le\beta$), on a les drapeaux ayant $\beta$ comme plus grand élément, de la forme $d\amalg d''_0$, où $d''_0\subset I_{>\alpha}$} et \struck{formé d'éléments} \add{des \uncertain{drapeaux}} $d''_0\subset I_{>i_p}$, contenant strictement $d'_0$ et tels que $d'_0$ soit cofinal dans $d''_0$, \add{\uncertain{l'ensemble de ces} $d''_0$, ou des $d''_0$, est en} \struck{\ill{}} correspondance bijective avec
\[
\mathrm{Drap}\big(I_{i_p<\cdot<i_{p+1}}\amalg I_{i_{p+1}<\cdot<i_{p+2}}\amalg\cdots\amalg I_{i_{n-1}<\cdot<i_n}\big)
\]
\note{le passage biffé, depuis « plus grand élément », est en partie encadré et barré d'une boucle ; au-dessus du symbole de la première flèche vers $\Sigma_{d'}$, un signe barré.}

\page{47}
\note{Suite de la page 46, sans numéro de sa main.}

On a lissité (\ill{} \ill{} \ill{}) ssi \struck{$i_p$} la suite $i_p\,i_{p+1}\ldots i_n$ (à partir de $i_p$, pas de $i_1$ !) est stricte, i.e.\ chaque terme est \add{(\uncertain{le plus proche}, au sens de $I$)} successeur du précédent.

\emph{11. Digression : \ill{} \uncertain{cas} \ill{} \ill{}} $I$.\note{titre souligné, lecture incertaine ; le paragraphe traite d'un ensemble ordonné topologique $\Sigma$ donné sans ensemble d'indices $I$.}

Soit donné un ensemble \uncertain{ordonné} topologique, i.e.\ un \struck{espace topologique} ens.\ ordonné $\Sigma=\Sigma(0),\Sigma(1)$, où $\Sigma(0)$ est un espace topologique, $\Sigma(1)$ une partie fermée de $\Sigma(0)\times\Sigma(0)$, qui soit le graphe d'une relation d'ordre. On suppose\marginal{\ill{} \uncertain{autres} \ill{}}

(a)\note{« a » cerclé.} $\Sigma(1)\xrightarrow{\ \mathrm{but}\ }\Sigma(0)$ est une immersion localement bordante, et propre (\uncertain{à fibres finies})

(b)\note{« b » cerclé, surchargé.} Cela implique que $\Sigma(0)\xrightarrow{\ \mathrm{diag}\ }\Sigma(1)$ est un \uncertain{isomorphisme} sur une partie ouverte et fermée. Soit $\Sigma^{*}(1)$\note{l'exposant, une petite étoile ou un point épais, est lu $*$ ; il distingue ces espaces de drapeaux stricts de $\Sigma(0)$, $\Sigma(1)$, et de l'indice $\Sigma_{*}$ des pages précédentes.} le complémentaire de l'image, i.e.\ l'ens.\ des inégalités strictes $i<j$. On désigne par $\Sigma^{*}(n)$ l'espace des drapeaux \emph{stricts} $i_0<i_1<\cdots<i_n$ de longueur $n$. Les $\Sigma^{*}(n)$ \struck{\ill{}} \uncertain{forment} un espace \add{pré}\struck{semi}-simplicial (pas d'applications de dégénérescence, \struck{\ill{}} ils sont dans les $\Sigma(n)$\ldots). \struck{Les} Si on a un

\page{48}
\note{En haut à gauche, « 20 » de sa main.}

diagramme d'inclusions \add{(croissantes)} d'ens.\ tot.\ ordonnés

(31)

\begin{tikzcd}[column sep=small]
 & d''=d\cup d' & \\
d \arrow[ur, hook] & & d' \arrow[ul, hook'] \\
 & \delta=d\cap d' \arrow[ul, hook] \arrow[ur, hook'] &
\end{tikzcd}

tel que ce soit un diagramme cocart.\ dans la cat.\ des ens.\ ordonnés (\uncertain{i.e.} $\delta$ segment initial de $d$, et \uncertain{contenant} l'élément final de $d$) alors le diagramme correspondant

(32)

\begin{tikzcd}[column sep=small]
 & \Sigma^{*}(d'') \arrow[dl] \arrow[dr] & \\
\Sigma^{*}(d) \arrow[dr] & & \Sigma^{*}(d') \arrow[dl] \\
 & \Sigma^{*}(\delta) &
\end{tikzcd}

est cartésien. \struck{(\uncertain{pour} \ill{} le cas \ill{})}\note{sur le schéma (31), les traits d'inclusion sont tracés sans tête nette ; on les oriente du plus petit au plus grand ensemble.}

\marginal{Plus précisément, on veut que pour $x\in\Sigma^{*}(0)$, les $z_\alpha\in\Sigma^{*}(1)$ au-dessus de $x$ (par but) sont en \emph{\uncertain{correspondance biunivoque}} avec les (b) germes $B_\alpha$ de sous-espaces \uncertain{bordants}, images \uncertain{locales} par \ill{} des germes \ill{} de $\Sigma^{*}(1)$ en les $z_\alpha$, \ill{} \ill{} \ill{} \ill{} \uncertain{points} \ill{} de $\Sigma^{*}(1)$ \ill{}, \uncertain{transversaux}, et \ill{} \ill{} correspondants, \ill{} \uncertain{distincts} \ill{} au-dessus de \ill{} \ill{} \ill{} les $z_\alpha$\ldots}\note{longue note en travers dans la marge gauche, depuis le milieu de la page jusqu'en bas, « b » cerclé ; lecture très incertaine hors des formules.}

Les \struck{\ill{}} \uncertain{conditions} b c d des n\textsuperscript{os} précédents \uncertain{s'écrivent} ici en une seule (du type d).

Considérons un ens.\ \struck{tot.} ordonné \add{de « dim » $\ge 1$} $d=(i_0,\ldots,i_n)$, \struck{$\Sigma^{*}(d)\to\Sigma(I)$} et considérons le diagramme

(33)\note{numéro écrit à gauche du diagramme, contre la note de la marge.}

\begin{tikzcd}[column sep=normal, row sep=small]
 & \Sigma^{*}(i_0,i_n) \arrow[dr, "\varphi_0"] & \\
\Sigma^{*}(d) \arrow[ur] \arrow[r] \arrow[dr] & \Sigma^{*}(i_1,i_n) \arrow[r, "\varphi_1"] & \Sigma^{*}(i_n) \\
 & \Sigma^{*}(i_{n-1},i_n) \arrow[ur, "\varphi_{n-1}"'] &
\end{tikzcd}

\note{entre $\Sigma^{*}(i_1,i_n)$ et $\Sigma^{*}(i_{n-1},i_n)$, des points de suspension verticaux.}

\struck{des applications bordantes} d'immersions, où $\varphi_0,\varphi_1,\ldots,\varphi_{n-1}$ sont bordantes. On suppose

\page{49}
\note{Suite de la page 48, sans numéro de sa main.}

que \struck{la famille est} \add{\uncertain{au voisinage}} de chaque \uncertain{pt} de $\Sigma^{*}(d)\to\Sigma^{*}(i_n)$ \add{dans $\Sigma^{*}(\lbrace i_n\rbrace)\simeq\Sigma^{*}(0)$}, $\Sigma^{*}(d)$ soit l'intersection des $\Sigma^{*}(\lbrace\alpha,i_n\rbrace)$ ($\simeq\Sigma^{*}(1)$), et que cette famille soit transversale.

Cela implique la chose analogue pour tout morphisme d'immersion
\[
\Sigma^{*}(d')\longrightarrow\Sigma^{*}(d)\qquad d\subset d',\ d,d'\ \text{cofinaux (avec le plus grand élément)}
\]
\struck{la structure}, on considère des $d''$ entre $d$ et $d'$ (donc cofinaux avec $d$) qui sont de la forme $d\cup\lbrace i\rbrace$ ($i\in d'\setminus d$) on regarde le diagramme

(34)

\begin{tikzcd}[column sep=normal, row sep=small]
 & \Sigma^{*}(d\cup j_1) \arrow[dr, "\text{bordant}"] & \\
\Sigma^{*}(d') \arrow[ur] \arrow[dr] & & \Sigma^{*}(d) \\
 & \Sigma^{*}(d\cup j_p) \arrow[ur, "\text{bordant}"'] &
\end{tikzcd}

\note{entre $\Sigma^{*}(d\cup j_1)$ et $\Sigma^{*}(d\cup j_p)$, des points de suspension verticaux ; les éléments de $d'\setminus d$ sont notés $j_1,\ldots,j_p$.}

qui généralise (33), et on \struck{suppose} veut qu'il satisfasse la même condition de transversalité. Cela \uncertain{signifie} que, les espaces envisagés, au voisinage des pts en jeu, i.e.\ des sous-espaces de
\[
\Sigma^{*}(\lbrace\alpha\rbrace)\simeq\Sigma^{*}(0),
\]
décrits \add{(grâce à b)} comme des intersections partielles des certains sous-espaces bordants de $\Sigma^{*}(\lbrace\alpha\rbrace)$, se coupant transversalement\ldots

\page{50}
\note{En haut à gauche, « 21 » de sa main.}

Il y a de plus des « conditions de circonstance » :

(c)\note{« c » cerclé.} L'application $\Sigma^{*}(1)\xrightarrow{\ \mathrm{source}\ }\Sigma^{*}(0)$ est \add{propre et} fibrante (« condition \add{d'}équisingularité »), cela implique que si $d\subset d'$ est une inclusion d'ens.\ finis tot.\ ordonnés, telle que $d$ soit \emph{segment initial} de $d'$, alors $\Sigma^{*}(d')\to\Sigma^{*}(d)$ est fibrante.

(d)\note{« d » cerclé.} \struck{Les} $\Sigma^{*}(0)$ est \struck{\ill{}} une variété à bord, \emph{et} $\Sigma^{*}(1)\overset{\mathrm{bord}}{\hookrightarrow}\Sigma^{*}(0)$ a comme image son bord. (« condition non singulière ») donc aussi les $\Sigma^{*}(d)$ (car un sous-espace bordant d'une variété à bord est une variété à bord \struck{\ill{}} \uncertain{de dimension} $<$ bord, $=$ dim que celui-ci ??) \uncertain{ou} on exige, \struck{\ill{} que les conditions \ill{}} \add{que les $\Sigma^{*}(d)$ soient des} variétés à bord\ldots)

\marginal{\uncertain{Dans tous les cas} \ill{}, \ill{} $\Sigma^{*}(1)\xrightarrow{\mathrm{bord}}\Sigma^{*}(0)$ \ill{} une partie \ill{} \uncertain{bordante}\ldots}\note{note en travers dans la marge gauche, à hauteur de (d) ; lecture très incertaine.}

\marginal{il faudrait vérifier que le bord de $\Sigma^{*}(d)$ est \uncertain{réunion} des $\Sigma^{*}(d')$ \uncertain{pour} $d\subset d'$ \uncertain{cofinal} \ill{} (??) \struck{\ill{}}, si c'est vrai, ça \uncertain{risque} d'être délicat ! (d'après ce qu'on a vu \ill{} $\Sigma$ \ill{} en dessous de $\lbrace 0,1,2\rbrace$) \ldots\ \uncertain{n'est pas} \ill{} une correspondance \ill{} de $d$)\ldots}\note{note en travers dans la marge gauche, en face de (e), le « ?? » cerclé ; lecture incertaine.}

(\uncertain{condition} \add{\uncertain{mixte}} \struck{\ill{} \ill{}} équisingulière-lisse)

(e)\note{« e » cerclé, dans la marge.} \struck{\ill{}} Le morphisme \add{source}
\[
\Sigma^{*}(1)\longrightarrow\Sigma^{*}(0)
\]
(qui est propre et fibrant) est de plus \struck{(une fibration à fibres des variétés} \add{à fibres des} variétés à bord, et le fibré des bords est \struck{le \ill{} de \ill{} par \ill{}} l'image de $\Sigma^{*}(2)$ par $\Sigma^{*}(2)\to\Sigma^{*}(1)$ associé : $\lbrace 0,2\rbrace\hookrightarrow\lbrace 0,1,2\rbrace$

\page{51}
\note{Suite de la page 50, sans numéro de sa main.}

\struck{\ill{}} Je \uncertain{reviens} sur la condition b), \uncertain{que j'ai précisée} beaucoup, sans \uncertain{vraiment} utiliser \uncertain{pleinement} l'information \uncertain{qui était} dans la \uncertain{version} précédente.

b$_1$) Pour $x\in\Sigma(0)=\Sigma^{*}(0)$, les $z_\alpha\in\Sigma^{*}(1)$ au-dessus de $x$ (par l'application \emph{but} : $\Sigma^{*}(1)\xrightarrow{\ b\ }\Sigma^{*}(0)$) sont en correspondance biunivoque avec les germes \add{$B_\alpha$} en $x$ de sous-espaces bordants de $\Sigma^{*}(0)$, \struck{\ill{}} images par $b$ des germes de $\Sigma^{*}(1)$ en les $z_\alpha$.

b$_2$) \struck{Il existe un} Si $n=\mathrm{card}(\lbrace z_\alpha\rbrace)$, il existe un élément \struck{\ill{}} unique de $\Sigma^{*}(n)$, soit $\tilde x$, tel que les $z_\alpha$ soient les images de $\tilde x$ par les \add{$n$ applications}
\[
\Sigma^{*}(n)\ \overset{\varphi_{0,n}}{\underset{\varphi_{n-1,n}}{\rightrightarrows}}\ \Sigma^{*}(1)
\]
de $\widetilde\Sigma(n)$ dans $\widetilde\Sigma(1)$.\note{il écrit trois flèches superposées, $\varphi_{0,n}$, $\varphi_{1,n}$, \ldots, $\varphi_{n-1,n}$, de $\Sigma^{*}(n)$ vers $\Sigma^{*}(1)\ni z_\alpha$.}

(\uncertain{Il en résulte} qu'il y a une \emph{structure} (l'une des « structures de \ill{} \add{\uncertain{totale}} (\ill{} \add{\uncertain{sur}} \uncertain{dans} $\Sigma^{*}(0)=\Sigma$) \ill{} 1 » $B_\alpha$ (\uncertain{passant par} $x$.)

\struck{\ill{}} b$_3$) Pour tout \struck{$1\le d\le n$, \ill{}} $1\le d\le n$, considérons l'application
\[
\Sigma^{*}(d)\longrightarrow\Sigma^{*}(0)\ni x
\]
\add{\uncertain{induite par}} \struck{\ill{}} l'inclusion \struck{\ill{}} $\lbrace d\rbrace\hookrightarrow\lbrace 0,1,\ldots,d\rbrace$, associée : les pts $z$ de $\Sigma^{*}(d)$ au-dessus de $x$ sont en correspondance biunivoque avec l'ens.\ des parties \add{de card $d$} d'un \struck{\ill{}} $I$ des germes $B_\alpha$, ou encore l'ens.\ des parties \add{de card $d$} de l'ens.\ $I$ des $z_\alpha$, ou encore \uncertain{avec} $\mathcal{P}_d([0,\ldots,n-1])$, en associant : \uncertain{à une telle} partie $A$ l'application \add{strictement} croissante unique de $\Delta_d=\lbrace 0,\ldots,d\rbrace$ \uncertain{dans}

\page{52}
\note{En haut à gauche, « 22 » de sa main.}

$\Delta_n=\lbrace 0,\ldots,n\rbrace$ \struck{dont} \uncertain{ayant pour} l'image \uncertain{soit} $A\cup\lbrace n\rbrace$, et l'hom.\ simplicial correspondant
\[
\underset{\tilde x}{\Sigma^{*}(n)}\ \xrightarrow{\ \text{imm.}\ }\ \underset{z}{\Sigma^{*}(d)}\ \Big(\overset{\text{imm.}}{\dashrightarrow}\ \underset{x}{\Sigma^{*}(0)}\Big)
\]
et l'image de $\tilde x$ par celui-ci.\note{sous chaque espace, le point qu'il contient ($\tilde x$, $z$, $x$), relié par $\in$ renversé ; sous $\Sigma^{*}(d)$ un griffonnage biffé, et après $\Sigma^{*}(d)$ une flèche biffée.} De plus, pour $z$ dans \struck{\ill{}} $\Sigma^{*}(d)$ au-dessus de $x$, \struck{$Z$} le germe de sous-espace de $\Sigma^{*}(0)$ en $x$, image du germe de $\Sigma^{*}(d)$ en $z$, est l'intersection des germes $B_\alpha$ correspondant aux $z_\alpha\in A$.

On s'intéresse ici aux morphismes entre les $\Sigma^{*}(n)$ correspondant aux applications strictement croissantes d'ens.\ tot\textsuperscript{t} ordonnés non vides $d\hookrightarrow d'$ tel que $d$ soit cofinal dans $d'$.

La catégorie de ces objets et morphismes équivaut : celle des ens.\ tot.\ ordonnés $d$ (\uncertain{év.}\ vides), par le foncteur qui à un tel $d$ associe $d\amalg\lbrace\omega\rbrace$, où $\omega$ est un « dernier élément », et à tout $d'\hookrightarrow d$ str.\ croiss.\ le prolongement par $\omega\mapsto\omega$. On a donc \struck{\ill{}} un foncteur
\[
\Gamma : d\longmapsto\Sigma^{*}(d\amalg\lbrace\omega\rbrace)
\]
sur les ens.\ tot.\ ordonnés \struck{\ill{}} (\uncertain{év.}\ vides), dont la donnée équivaut : celle des $\Sigma^{*}(n)$, avec les \struck{\ill{}} appl.\ simpliciales « cofinales » entre eux. Et la donnée d'un tel foncteur $\Gamma$ revient : la donnée de $\Gamma(\emptyset)$ (ici $\Sigma^{*}(0)$) et

\page{53}
\note{Suite de la page 52, sans numéro de sa main.}

d'un objet (espace) simplicial \emph{au-dessus de} $\Gamma(\emptyset)=\Sigma^{*}(0)$. Ici il s'agit d'un espace simplicial relatif, et les conditions b$_1$, b$_2$, b$_3$ concernent cette structure -- elles disent \uncertain{intuitivement} que cette structure est associée : une stratification cylindrique locale de $\Sigma^{*}(0)=\Sigma$, et « ordonnée » (i.e.\ les strates de codim.\ 1 passant par un $x\in\Sigma$ sont ordonnées, et cette relation d'ordre varie continûment avec $x$ en un sens évident).\marginal{NB $\Gamma(n)=\Sigma^{*}(n+1)$}

Les \struck{$\Sigma^{*}(d)$ s'id.} $\Gamma(d)$ \add{$=\Sigma^{*}(d+1)$} s'identifient alors aux \struck{parties} espaces des \add{(germes de)} strates de codim.\ $d+1$, \struck{les} \add{germes} identifiés \struck{(grâce à la structure d'ordre)} aux ens.\ de germes de strates de codim.\ 1 \struck{\ill{}} $\ni x$. Les applications semi-simpliciales entre les $\Gamma(d)$ proviennent alors de la structure d'ordre -- \struck{\ill{}} par \uncertain{reconstitution} les espaces de drapeaux associés : la stratification, à l'aide des \struck{\ill{}} $\Gamma(d)$, grâce : cette structure ordonnée. En résumé, la donnée de $\Gamma : d\mapsto\Gamma(d)$ \struck{\ill{}} (i.e.\ de $\Sigma^{*}$ avec les \uncertain{app.}\ semi-simpliciales cofinales) équivaut : celle d'un $\Sigma=\Sigma^{*}(0)$,

\page{54}
\note{En haut à gauche, « 23 » de sa main.}

avec la structure de stratification cylindrique \add{ordonnée} locale. \struck{Par exemple,} Une autre façon de décrire les $\Gamma(d)=\Sigma^{*}(d+1)$, c'est de prendre la puiss.\ \struck{sym.}\ \add{cart.\ relative} $d$-ième de $\Gamma_0$ \add{$=\Sigma^{*}(1)$} \struck{\ill{}} sur $\Sigma$ ($=\Gamma(\emptyset)$), \uncertain{d'où} on \uncertain{utilise} les images \add{$\Gamma'(d)$} en dim.\ ($d'\le d$) des puiss.\ cart.\ relatives \struck{définies} par les applications ½ simpliciales entre celles-ci. Cela donne : \uncertain{préciser} les \struck{\ill{}} suites ordonnées de $d$ germes de strates de codim.\ 1 \add{\uncertain{en un pt}} (\struck{\ill{}} distinctes), soit $\Gamma'(d)$. On a alors $\Gamma(d)\subset\Gamma'(d)$, $\Gamma(d)$ est une partie : la fois ouverte et fermée de $\Gamma'(d)$, et \struck{$\Gamma'(d)$ \ill{}}
\[
(35)\qquad \Gamma'(d)\simeq\mathfrak{S}_d\times\Gamma(d)
\]
\note{« ½ simpliciales » : le signe ½ pour \emph{semi-}, comme ailleurs dans le fonds. Au-dessus de $\Gamma(d)\subset\Gamma'(d)$, un exposant peu net sur le premier $\Gamma$.}

Pour avoir \struck{la} \add{une} structure d'ordre \add{totale} relative, de \struck{\ill{}} $\Gamma_0=\Sigma^{*}(1)$ sur $\Sigma=\Sigma^{*}(0)$, il \add{\uncertain{suffit}} \struck{suffit} de se donner, \struck{dans $\Gamma_1(2)=\Sigma^{*}(2)$, $\Gamma'$ (\ill{})} dans $\Gamma_0\times_\Sigma\Gamma_0\setminus\text{diagonale}$, une partie (ouverte et fermée) $\Gamma_1$, telle que
\[
\Gamma_1\cup{}^{s}\Gamma_1=\Gamma_0\times\Gamma_0\setminus\mathrm{diag},\qquad \Gamma_1\cap{}^{s}\Gamma_1=\emptyset
\]
\uncertain{i.e.}\ une section de $\Gamma_0\times_\Sigma\Gamma_0\setminus\mathrm{diag}$ au-dessus de $(\mathrm{Sym}^2_\Sigma\Gamma_0)^{*}$ (« choix d'un premier élément parmi deux éléments distincts »), avec la condition de transitivité
\[
(36)\qquad (\Gamma_1\times_\Sigma\Gamma_0)\cap(\Gamma_0\times_\Sigma\Gamma_1)\xrightarrow{\ \mathrm{pr}_{13}\ }\Gamma_0\times\Gamma_0\quad\text{se factorise par }\Gamma_1,
\]
où l'intersection est prise dans $\Gamma_0\times\Gamma_0\times\Gamma_0$ et $\Gamma_1\subset\Gamma_0\times\Gamma_0$.\note{${}^{s}\Gamma_1$ : le symétrique de $\Gamma_1$. Les deux inclusions de la dernière ligne sont écrites sous la formule, verticalement ; « pr$_{13}$ » surcharge un premier indice. Devant (36), un mot biffé.}

\marginal{\struck{il s'agit de} \ldots\ $\Gamma_1$ \uncertain{se voit} \ill{} $x<y$ \ldots\ \uncertain{relations} \ldots\ \struck{ou encore une section \ill{} de $\Gamma_0\times_\Sigma\Gamma_0$ \ill{} de $\mathrm{Sym}^2(\ldots)$}}\note{notes en travers dans la marge gauche, en partie biffées ; lecture très incertaine.}

\page{55}
\note{Sans numéro de sa main ; suite du feuillet 23 (page 54).}

Je m'aperçois finalement que j'ai pas mal bafouillé pour arriver : trouver les formulations simples -- car j'ai \emph{mélangé} des propriétés purement ensemblistes de l'ens.\ ordonné $\Sigma=\Sigma^{*}(0)$, \struck{\ill{}} ses interprétations semi-simpliciales et \struck{\ill{}} \add{les} aspects topologiques. Je \emph{reprends}

\struck{$\Sigma$ ens.}

$\Sigma=\Sigma(0)$ ens.\ ordonné, muni d'une topologie \struck{séparée} (\uncertain{pas} \ill{}

Si $\Sigma_1$ est le graphe de la relation d'ordre, on suppose

(a)\note{« a » cerclé.} $\Sigma(1)\xrightarrow{\ b=\mathrm{pr}_2\ }\Sigma(0)$ est propre et loc.\ immersif bordant (propriétés topologiques) \struck{\ill{}}

\struck{(b)}\note{la lettre cerclée est biffée.} Cela implique \struck{par l'implication} que l'application diagonale est ouverte, donc si $\Sigma$ séparé, son image est une partie : la fois ouverte et fermée -- et on peut alors introduire
\[
\Sigma^{*}_{f}(1)=\Sigma(1)\setminus\text{diagonale}.
\]
Si $\Sigma$ est séparé, la propriété de $\Sigma(1)\to\Sigma(0)$ implique aussi que $\Sigma_1$ est \emph{fermé} dans $\Sigma\times\Sigma$ \struck{$\Sigma(0)\times\Sigma(0)$}.\note{l'indice $f$ de $\Sigma^{*}_{f}(1)$ est une lecture douteuse.}

(b)\note{« b » cerclé.} Pour tout $x\in\Sigma$, l'ens.\ \add{(fini)} $\Sigma_{\le x}$ \add{(fibre de $x$ dans $\Sigma(1)$)} est \emph{tot.\ ordonné} par l'ordre induit.

\struck{(c)} Donc la donnée d'un drapeau \add{$x_{i_0}\ldots x_{i_d}$} \struck{de longueur} $d$ dans $\Sigma_{\le}$ équivaut : celle d'une partie de card.\ $d+1$

\page{56}
\note{En haut à gauche, « 24 » de sa main.}

de l'ens.\ à $N(x)$ éléments, où $N(x)=\mathrm{card}(\Sigma_{\le x})$. On \uncertain{a encore}, si on regarde l'application-but
\[
\Sigma^{*}(d)\longrightarrow\Sigma=\Sigma^{*}(0),\qquad (x_0<\cdots<x_d)\longmapsto x_d
\]
la fibre de $x\in\Sigma$ s'identifie : l'ens.\ \struck{des parties} $\mathcal{P}_d(\lbrace 1,\ldots,N(x)\rbrace)$.

La structure relative de $\Sigma^{*}(1)$ sur $\Sigma=\Sigma^{*}(0)$ est donc celle d'un ens.\ tot.\ ordonné relatif, qui permet de reconstituer le système des $\Sigma^{*}(d)$ ($d\in\mathbb{N}$) avec les applications semi-simpliciales \add{(i.e.\ : \uncertain{toute condition})} correspondant aux morphismes cofinaux entre ens.\ tot.\ ordonnés-types.\marginal{tout ceci est purement ensembliste (\ill{}) indépendant des top.\ \ill{}}\note{note en travers dans la marge gauche, avec un petit rond cerclé.}

(c)\note{« c » cerclé.} $\forall x\in\Sigma$, les germes \add{$B_\alpha$} de sous-espaces \add{bordants} de $\Sigma$ en $x$ correspondants, par (a), aux $z_\alpha\in\Sigma^{*}(1)$ au-dessus de $x$ (par $\Sigma^{*}(1)\xrightarrow{\ b\ }\Sigma$) sont mutuellement transversaux (et deux : deux distincts).\note{le renvoi « par (a) » est un petit rond cerclé, lu comme (a).}

Les germes intersections de certains des $B_\alpha$ correspondent donc aux pts des $\Sigma^{*}(d)$ au-dessus de $x$ (par $\Sigma^{*}(d)\xrightarrow{\ b\ }\Sigma$), comme images du \struck{\ill{}} germe de $\Sigma^{*}(d)$ en ces pts.

\struck{\ill{}} Les espaces $\Sigma^{*}(d)$ ($d\in\mathbb{N}$) et les applications ½ simpliciales cofinales entre eux, \struck{\ill{}} \uncertain{ou encore} l'espace relatif $\Sigma^{*}(1)$ sur $\Sigma^{*}(0)$ et la structure d'ordre relative, \uncertain{ces} décrivent entièrement

\page{57}
\note{Sans numéro de sa main ; suite du feuillet 24 (page 56).}

\struck{\ill{}} $\Sigma$, muni d'une stratification locale cylindrique ordonnée : $\Sigma^{*}(1)$ est l'espace des strates de codim.\ 1, $\Sigma^{*}(1)\xrightarrow{\ b\ }\Sigma$ l'application associée : un germe sa \uncertain{trace}, $\Sigma^{*}(2)$ est l'espace des paires de germes $B,B'$ en $x$ \struck{\ill{}} tels que $\lbrace B\rbrace<\lbrace B'\rbrace$\ldots

\note{Dans la marge gauche, en face, une colonne ; entre $\Sigma^{*}(2)$ et $\Sigma^{*}(1)$ il trace deux flèches parallèles, $s'$ et $b'$, réunies ici en une seule.}

\begin{tikzcd}[row sep=small]
\vdots \\
\Sigma^{*}(2) \arrow[d, "{s',\ b'}"] \\
\Sigma^{*}(1) \arrow[d, "b"] \\
\Sigma=\Sigma^{*}(0)
\end{tikzcd}

Pour retrouver la relation d'ordre sur $\Sigma=\Sigma^{*}(0)$ \add{à partir de la structure d'ordre relative $\Sigma^{*}(1)$ sur $\Sigma^{*}(0)$ (ou de ce que $\Sigma^{*}(1)$ \struck{\ill{}} \uncertain{soit} \struck{\ill{}} \add{\uncertain{\ill{}}} les relations strictes $x<y$ \uncertain{dans} $\Sigma$)} \struck{\ill{}} on doit se donner une application \struck{\ill{}} \uncertain{source}
\[
(36)\qquad \Sigma^{*}(1)\xrightarrow{\ \underline s=\mathrm{pr}_1\ }\Sigma=\Sigma^{*}(0)
\]
qui soit injective \struck{et croissante} sur chaque fibre \add{$\underline b^{-1}(x)$} de $\Sigma^{*}(1)\xrightarrow{\ \underline b\ }\Sigma$ (i.e.\ $\Sigma^{*}(1)\xrightarrow{(\underline s,\underline b)}\Sigma\times\Sigma$ inj.) et d'image $\not\ni x$ (i.e.\ $\mathrm{Im}(\underline s,\underline b)\cap\mathrm{diag}(\Sigma\times\Sigma)=\emptyset$)\note{il écrit d'abord « $\Sigma_\emptyset$ » (surchargé) pour la base ; le numéro (36) est déjà celui de la condition de transitivité de la page 54 : transcrit tel quel.}

\struck{et on doit se donner de plus (37) $\Sigma^{*}(2)\xrightarrow{\ \varphi_2\ }\Sigma^{*}(1)\ (\hookrightarrow\Sigma\times\Sigma)$} \struck{\ldots\ (NB la donnée de $\varphi_2$ équivaut : celle des deux applications $\mathrm{pr}_1\circ\varphi,\ \mathrm{pr}_2\circ\varphi : \Sigma^{*}(2)\rightrightarrows\Sigma$ qu'on se donne\ldots)} \struck{(38) $\Sigma^{*}(2)\xrightarrow{(\mathrm{pr}_1\circ\varphi,\ \mathrm{pr}_2\circ\varphi,\ \psi)}\Sigma_0\times\Sigma_0\times\Sigma_0$} \struck{dont l'image est par construction (voir (37)) dans $\Sigma_1\times\Sigma_0$ et dans\ldots}\note{Toute la moitié inférieure de la page, depuis « et on doit se donner de plus », est biffée de quatre longs traits obliques et s'arrête en cours de phrase ; le bas de la page est blanc. Dans la marge gauche, également sous les traits : « il \ill{} \uncertain{de} \ill{} \ill{} $\mathrm{pr}_2$ \ill{} \ill{} \ill{} $\mathrm{pr}_1$ \ill{} $\mathrm{pr}_{13}$ \ill{} ». Le « (37) » du renvoi est cerclé.}

\page{58}
\note{En haut à gauche, « 25 » de sa main, entouré d'accolades.}

de telle façon qu'il existe
\[
(37)\qquad p_2 : \Sigma^{*}(2)\longrightarrow\Sigma^{*}(1)
\]
tel que
\[
(38)\qquad \underline b\circ p_2=\underline s\circ\underline b'\quad\text{i.e.\ le carré}
\]

\begin{tikzcd}
\Sigma^{*}(2) \arrow[r, hook, "\underline b'"] \arrow[d, "p_2"'] & \Sigma^{*}(1) \arrow[d, "p_1=\underline s"] \\
\Sigma^{*}(1) \arrow[r, hook, "\underline b"'] & \Sigma
\end{tikzcd}

\note{au centre du carré, « (cart.) » ; les flèches horizontales partent d'un petit crochet en pointillé. La phrase « est commutatif », écrite plus bas, vaut pour les deux carrés.}

\[
p_1\circ p_2=p_1\circ\underline s'\quad\text{i.e.\ le carré}
\]

\begin{tikzcd}[column sep=small, row sep=small]
 & \Sigma^{*}(1) \arrow[dl, "p_1"'] & \\
\Sigma & & \Sigma^{*}(2) \arrow[ul, "\underline s'"'] \arrow[dl, "p_2"] \\
 & \Sigma^{*}(1) \arrow[ul, "p_1"] &
\end{tikzcd}

est commutatif.\note{les deux égalités de (38) sont réunies par une accolade ; entre les deux, une ligne en travers, « \ill{} \uncertain{relatif} », séparée par un trait.}

Ce qui détermine $p_2$ de façon unique, puisque $(\struck{\ill{}}\,p_1,\underline b) : \Sigma^{*}(1)\to\Sigma\times\Sigma$ est mono).

De plus on \struck{suppose} \add{doit} exiger que le premier carré (38) soit \emph{cartésien}, et que $p_1 : \Sigma^{*}(1)\to\Sigma$ soit \emph{transversal} p.r.\ à l'immersion bordante \add{\uncertain{l'unique}} $\underline b : \Sigma^{*}(1)\overset{\underline b}{\hookrightarrow}\Sigma$ (ce qui est \struck{\ill{}} \add{une} condition de nature non purement algébrique\ldots)

Posons alors, pour $x,y\in\Sigma$
\[
(39)\qquad x\prec y\iff\exists\,u\in\Sigma^{*}(1)\ \text{tel que}\ x=p_1u,\ y=\underline b u
\]

\page{59}
\note{Sans numéro de sa main ; suite du feuillet 25 (page 58).}

On voit alors que
\[
x\prec y\ \text{et}\ y\prec z\ \Longrightarrow\ x\prec z
\]
car soient $u,v\in\Sigma^{*}(1)$ tels que
\[
p_1u=x,\quad \underline b u=y,\quad p_1v=y,\quad \underline b v=z
\]
on a donc
\[
\underline b u=p_1v
\]
donc par le carré cartésien de (38), \struck{$\exists w$}
\[
\exists!\,d\in\Sigma^{*}(2),\ \text{avec}\quad \underline b'(d)=v,\quad p_2(d)=u
\]
soit
\[
w=\underline s'(d)\in\Sigma^{*}(1)
\]
on a
\[
\begin{aligned}
p_1(w)&=p_1(\underline s'(d))\overset{(38)}{=}p_1p_2(d)=p_1(u)=x\\
\underline b(w)&=\underline b(\underline s'(d))=\underline b(\underline b'(d))=\underline b(v)=z
\end{aligned}
\]
donc
\[
w : x\prec z .
\]
\note{sous « $\underline b u=y,\ p_1v=y$ », une accolade réunit les deux égalités. Dans la seconde ligne, $\underline b(\underline s'(d))=\underline b(\underline b'(d))$ est écrit tel quel.}

De plus, on a
\[
x\not\prec x\qquad\big(\text{car } \mathrm{Im}\big(\Sigma^{*}(1)\xrightarrow{(p_1,\underline b)}\Sigma\times\Sigma\big)\cap\mathrm{diag}=\emptyset\big)
\]
donc la relation $x\prec y$ est la relation d'inégalité stricte associée : une relation d'ordre sur $\Sigma$. On récupère $\Sigma^{*}(1)$ comme le graphe de la relation d'ordre stricte, $\Sigma^{*}(2)$ comme les drapeaux $x\prec y\prec z$ de longueur 2 dans $\Sigma$, et les applications

\page{60}
\note{En haut à gauche, « 26 » de sa main.}

\[
\Sigma^{*}(1)\ \overset{p_1=\underline s}{\underset{\underline b}{\rightrightarrows}}\ \Sigma\qquad\qquad \Sigma^{*}(2)\ \overset{p_2}{\underset{\underline b'}{\rightrightarrows}}\ \Sigma^{*}(1)
\]
\note{à droite, trois flèches parallèles, étiquetées $p_2$, $\underline s'$ et $\underline b'$ ; on n'en a gardé que les deux extrêmes dans la formule.}

comme
\[
(39)\qquad\left\lbrace
\begin{array}{l}
(x\prec y)\xrightarrow{\ p_1=\underline s\ }x\\
(x\prec y)\xrightarrow{\ \underline b\ }y\\
(x\prec y\prec z)\xrightarrow{\ p_2\ }(x,y)\\
(x\prec y\prec z)\xrightarrow{\ \underline s'\ }(x,z)\\
(x\prec y\prec z)\xrightarrow{\ \underline b'\ }(y,z)
\end{array}\right.
\]
\note{le numéro (39) est déjà celui de la définition de $\prec$, page 58 : transcrit tel quel. Dans les deux dernières lignes, le premier terme de la paire est surchargé ($x$, resp.\ $y$, écrits par-dessus une autre lettre).}

\emph{Scholie} En résumé, la situation est décrite par ces données \struck{suivantes}

1°) Un espace top.\ \struck{\uncertain{séparé}} $\Sigma$, muni d'une stratification (pré)cylindrique -- soit $\Sigma_1$ l'espace des strates de codim.\ 1, de sorte que $\Sigma^{*}(1)\xrightarrow{\ \underline b\ }\Sigma$ est une imm.\ bordante \add{propre}.

2°) Une relation d'ordre totale relative de $\Sigma_1$ sur $\Sigma$ (définie par la donnée, pour tout $x\in\underline b(\Sigma_1)$, d'un ordre total sur l'ens.\ des strates de codim.\ 1 qui passent par $x$, cette relation d'ordre totale étant soumise : la condition d'être \add{donnée par} une partie \emph{fermée} de $(\Sigma_1\times_\Sigma\Sigma_1)\ldots$) -- soit
\[
\Sigma^{*}(2)\ \overset{\underline s'}{\underset{\underline b'}{\rightrightarrows}}\ \Sigma^{*}(1)\xrightarrow{\ \underline b\ }\Sigma
\]
\note{sous « $(\Sigma_1\times_\Sigma\Sigma_1)$ », $\Sigma^{*}(2)$ est écrit à gauche, en renvoi ; les indices de ce produit sont surchargés.}

\marginal{« stratification cylindrique loc.\ tot.\ ordonnée »}\note{note en travers dans la marge gauche, reliée par une accolade à 1°) et 2°) ; lecture incertaine.}

3°) Une application continue
\[
p_1=\underline s : \Sigma^{*}(1)\longrightarrow\Sigma
\]
ayant les propriétés suivantes

a) $(\underline s,\underline b) : \Sigma^{*}(1)\longrightarrow\Sigma\times\Sigma$ est injectif, et son image ne rencontre pas la diagonale, i.e.\ $\forall u\in\Sigma^{*}(1)$, $\underline s(u)\neq\underline b(u)$

\end{document}
